\section{Derivations}
\label{app:derivations}

This appendix carries the mathematics behind the study: the value of a
concentrated position, the two loss measures, the concentration factor, the
fee and hedge accounting, the risk measure and the policy update. The results
in the main body use these expressions as definitions, so the appendix doubles
as a statement of exactly what was measured.

\subsection{The value of a concentrated position}

A constant-product pool with reserves $(x,y)$ quotes the invariant
$xy=k$ and a price $P=y/x$. A Uniswap~v3 position replaces the full-range
reserves with a liquidity constant $L$ attached to a band
$[p_\ell,p_u]$ in price space; the band's boundaries are tick indices
$i_\ell,i_u$, which correspond monotonically to $p_\ell,p_u$ through the
pool's conversion between ticks and prices. The position's token amounts follow from the
invariant in square-root coordinates. Writing $\sqrt{P}$ for the current
square-root price,

\begin{align}
x(P) &= L\big(\sqrt{P}-\sqrt{p_\ell}\big),
\label{eq:app-x}\\
y(P) &= L\Big(\frac{1}{\sqrt{P}}-\frac{1}{\sqrt{p_u}}\Big),
\label{eq:app-y}
\end{align}
so $x$ is the USDC amount and $y$ the WETH amount, matching the main body's
token0/token1 convention. While $p_\ell\le P\le p_u$; outside the band the
position holds only the token the price has moved toward. The mark value in
USDC is $V(P)=x(P)+y(P)P$, which in range becomes

\begin{equation}
V(P) = L\Big(2\sqrt{P}-\sqrt{p_\ell}-\frac{P}{\sqrt{p_u}}\Big).
\label{eq:app-value}
\end{equation}

The deposit is fixed at $W_0$ and the position is opened at the entry price
$P_0$ with the current tick as the band's centre, so the liquidity is pinned
by $V(P_0)=W_0$:

\begin{equation}
L = W_0 \Big/ \Big(2\sqrt{P_0}-\sqrt{p_\ell}-\frac{P_0}{\sqrt{p_u}}\Big).
\label{eq:app-L}
\end{equation}

The position's exposure is the derivative of the value with respect to the
price. In range,

\begin{equation}
\frac{\partial V}{\partial P} = y(P) = L\Big(\frac{1}{\sqrt{P}}-\frac{1}{\sqrt{p_u}}\Big),
\qquad
\frac{\partial^2 V}{\partial P^2} = -\frac{L}{2}\,P^{-3/2}<0.
\label{eq:app-greeks}
\end{equation}

The LP is long the risky token while the price is low in the band, and short
it while the price is high; the value is concave in the price because the
position sells the asset that rises and buys the one that falls. Concavity is
the source of the loss to arbitrage: no matter which way the price moves, the
position underperforms a portfolio that would have rebalanced to the same
exposure continuously.

\subsection{Impermanent loss}

The HODL benchmark holds the entry amounts $(x_0,y_0)=(x(P_0),y(P_0))$ fixed,
so its value is $H(P)=x_0+y_0P$. Impermanent loss is the gap

\begin{equation}
\mathrm{IL}(P) = V(P)-H(P)\;\le\;0\quad\text{for }P\ne P_0,
\label{eq:app-il}
\end{equation}
with equality at entry. Substituting \eqref{eq:app-value} and the entry
amounts gives the closed form used by the evaluation:

\begin{equation}
\frac{V(P)}{W_0} = \frac{2\sqrt{P}-\sqrt{p_\ell}-P/\sqrt{p_u}}
{2\sqrt{P_0}-\sqrt{p_\ell}-P_0/\sqrt{p_u}},
\qquad
\frac{H(P)}{W_0} = \frac{x_0}{W_0}+\frac{y_0}{W_0}P.
\label{eq:app-il-closed}
\end{equation}

For a band symmetric in ticks around $P_0$, write $p_\ell=P_0/r$ and
$p_u=P_0 r$ with $r=e^{\delta}$, where $\delta$ is the half-width in
log-price and equals the half-width in ticks multiplied by the tick size. At
the entry price the two legs are equal in value,

\begin{equation}
\frac{x_0}{W_0} = \frac{y_0P_0}{W_0} = \frac{1}{2},
\qquad
\frac{y_0}{W_0} = \frac{1}{2P_0},
\label{eq:app-amounts}
\end{equation}
so the deposit is balanced at any symmetric width and
$x_0/W_0+y_0P_0/W_0=1$ by construction.

\subsection{Concentration and loss-versus-rebalancing}

Consider a price that follows a geometric Brownian motion,
$\mathrm{d}P_t=\mu P_t\,\mathrm{d}t+\sigma P_t\,\mathrm{d}W_t$. Let
$R_t$ be the value of the continuously rebalanced portfolio that starts at
$V(P_0)$ and holds the AMM's instantaneous exposure at every instant,
$\mathrm{d}R_t=({\partial V}/{\partial P})\,\mathrm{d}P_t$. It\^o's lemma
gives

\begin{equation}
\mathrm{d}V_t = \frac{\partial V}{\partial P}\,\mathrm{d}P_t
+ \frac{1}{2}\frac{\partial^2 V}{\partial P^2}\,\sigma^2P_t^2\,\mathrm{d}t,
\label{eq:app-ito}
\end{equation}
so the gap to the rebalancing portfolio accumulates at

\begin{equation}
\mathrm{dLVR}_t = -\frac{1}{2}\frac{\partial^2 V}{\partial P^2}\,\sigma^2P_t^2\,\mathrm{d}t
\;\ge\;0 ,
\label{eq:app-lvr-general}
\end{equation}
which is the loss-versus-rebalancing rate of Milionis et al.\cite{milionis2022lvr}.
Substituting the curvature of \eqref{eq:app-greeks}, an in-range concentrated
position loses at the instantaneous rate

\begin{equation}
\mathrm{dLVR}_t = \frac{\sigma^2}{4}\,L\,\sqrt{P_t}\,\mathrm{d}t .
\label{eq:app-lvr}
\end{equation}

The full-range position is the special case $p_\ell\to0$, $p_u\to\infty$ with
$V=2L\sqrt{P}$, for which \eqref{eq:app-lvr} collapses to
$\sigma^2V/8$, the rate in the LVR literature. Dividing the two rates and
evaluating at the entry price, with both positions pinned to the same capital
$W_0$, gives the position's concentration factor,

\begin{equation}
\kappa = \frac{\mathrm{dLVR}^{\text{concentrated}}}
{\mathrm{dLVR}^{\text{full range}}}
= \frac{1}{1-(p_\ell/p_u)^{1/4}}
= \frac{1}{1-e^{-\delta/2}} ,
\label{eq:app-kappa}
\end{equation}
so a band of half-width $\delta=0.012$ in log-price, the default open band of
the replay (half-width $120$ ticks), carries an instantaneous LVR rate about
$167$ times the full-range
rate while it is in range. The factor is why a narrow band earns more fees and
loses more to arbitrage at the same time, and why the fee tier, not the width
alone, decides whether the position clears its own LVR.

Fees modify the arbitrage relation. With a proportional fee $\iota$, a
riskless arbitrageur only trades when the pool price deviates from the external
price by more than the fee, so the pool price wanders inside a band of width
proportional to $\iota$ around the external price before arbitrage is
profitable. Fee income and the residual LVR are then two sides of the same
threshold, which is the object the empirical fee-to-LVR ratio in
Appendix~\ref{app:ablations} measures; the fee-adjusted arbitrage condition is
derived by Milionis et al.\cite{milionis2023fees}.

\subsection{Fee income and the liquidity share}

While the price is inside the band, the position earns a proportional share of
every swap fee. If the position's liquidity is $L$ and the pool's active
liquidity is $L_{\text{active}}$, that share is the ratio $L/L_{\text{active}}$
on each crossed tick segment. The exact engine does not approximate this ratio
at the pool level: each swap is apportioned across the tick segments it
crosses, and the position's fee accrual is the integer change in the inside
fee-growth accumulator multiplied by its liquidity. Fee income over a step is
therefore

\begin{equation}
f_t = \sum_{\text{segments } s} L_s \cdot
\big(\Delta \text{feeGrowthInside}_s\big) \big/ 2^{128} ,
\label{eq:app-fee}
\end{equation}
with $L_s=0$ outside the band and the $2^{128}$ divisor the protocol's
fixed-point normalisation. The \emph{marginal-agent} assumption of the
simulator replaces the sum with the position's share of the global fee growth
while in range, which removes the segment bookkeeping at the cost of a small
approximation bounded by the parity harness.

\subsection{The hedge leg}

The hedge shorts $h_t\,y(P_t)$ units of WETH through a perpetual swap, where
$y$ is the LP's WETH amount. Its mark PnL over a step is

\begin{equation}
\pi_t = -\,h_{t-1}\,y(P_{t-1})\,\big(P_t-P_{t-1}\big),
\label{eq:app-hedge-pnl}
\end{equation}
funding is $\phi_t=+h_{t-1}y(P_{t-1})P_{t-1}\,\rho_t$ when the eight-hour
rate $\rho_t$, scaled to the step, is positive (the short collects it), and
adjusting the position costs a taker fee $\tau_t\le0$ proportional to the
change in short notional. The joint
strategy value is $V_t+\text{fees}_t+\pi_t+\phi_t+\tau_t-g_t-S_t$. Because the
hedge is linear in the price and the LP value is concave, the hedge removes
the first-order exposure but cannot remove the curvature term
\eqref{eq:app-lvr}: a delta-neutral LP still pays LVR. This is the structural
reason the empirical hedge cuts the tail without repairing the mean.

\subsection{Risk measure and the training objective}

For a loss $L=-G$ on the episode gain $G$, conditional value-at-risk at level
$\alpha$ is the mean of the worst $1-\alpha$ tail,

\begin{equation}
\mathrm{CVaR}_\alpha(L) = \frac{1}{1-\alpha}\int_\alpha^1
\mathrm{VaR}_u(L)\,\mathrm{d}u
= \min_{\nu\in\mathbb{R}}\Big\{\nu+\frac{1}{1-\alpha}\mathbb{E}\big[(L-\nu)^+\big]\Big\},
\label{eq:app-cvar}
\end{equation}
the convex representation of Rockafellar and Uryasev\cite{rockafellar2000cvar};
CVaR is a coherent risk measure in the sense of
Artzner et al.\cite{artzner1999coherent}. The reward subtracts a multiple of a
trailing empirical estimate of \eqref{eq:app-cvar} on per-step joint equity
changes, with a one-day window, so the policy's gradient sees the same tail
statistic the evaluation reports.

Policy optimisation uses the clipped surrogate of PPO\cite{schulman2017ppo}.
With probability ratio $r_t(\theta)$ between the new and old policies and
advantage $\hat A_t$ from generalised advantage estimation\cite{schulman2015gae},

\begin{equation}
\mathcal{L}(\theta) = \mathbb{E}\Big[
\min\big(r_t(\theta)\hat A_t,\;
\mathrm{clip}(r_t(\theta),1-\epsilon,1+\epsilon)\hat A_t\big)\Big]
- c_v\,\mathbb{E}\big[(V_\theta(s_t)-\hat V_t)^2\big]
+ c_e\,\mathbb{E}\big[H(\pi_\theta(\cdot\mid s_t))\big],
\label{eq:app-ppo}
\end{equation}
with clip $\epsilon=0.15$, value coefficient $c_v$ and entropy bonus $c_e$.
The advantage is the exponentially weighted sum of one-step temporal-difference
residuals with parameter $\lambda_{\mathrm{GAE}}=0.95$. The full training
settings are the subject of Appendix~\ref{app:system}.
